\documentclass[11pt,a4paper]{article}
\usepackage[utf8]{inputenc}
\usepackage[margin=22mm]{geometry}
\usepackage{amsmath,amssymb,amsfonts}
\usepackage{booktabs}
\usepackage[hidelinks]{hyperref}
\usepackage{fancyhdr}

\pagestyle{fancy}
\fancyhf{}
\rhead{\small CALPHADHEACalc Reference Manual}
\lhead{\small SRM University-AP | Department of Mechanical Engineering}
\cfoot{\thepage}
\renewcommand{\headrulewidth}{0.4pt}

\title{\textbf{CALPHADHEACalc: Governing Equations and Variable Reference}\\[0.3em]
\large Thermodynamic Formulations, Physical Models, and Implementation Notes}

\author{
  \textbf{Felix Mwiya} \quad \textbf{Prof. Dr. Sheela Singh}\\[0.4em]
  \small Department of Mechanical Engineering, School of Engineering and Sciences\\
  \small SRM University-AP, Andhra Pradesh 522240, India
}
\date{October 2026}

\begin{document}
\maketitle

\begin{abstract}
This reference manual documents the thermodynamic, mechanical, kinetic, and transport equations implemented in the CALPHADHEACalc software. For each model, we state the governing relationship, the physical meaning and dimensions of its variables, the corresponding software subroutines, and the supporting literature citations. The formulations cover CALPHAD solution thermodynamics, empirical solid-solution criteria, soft-sphere lattice distortion, yield strength superposition, transport properties across the solidus-liquidus window, high-temperature oxidation and electrochemical corrosion kinetics, and kinematic X-ray diffraction.
\end{abstract}

\tableofcontents
\vspace{1em}
\hrule
\vspace{1.5em}

\section{CALPHAD Solution Thermodynamics and Phase Equilibria}
Primary Modules: \texttt{HEAThermoEngine.m}, \texttt{HEATDBEngine.m}

\subsection{Molar Gibbs Free Energy of Phase $\phi$}
\begin{equation}
G_m^\phi(T, \mathbf{c}) = {^{\text{ref}}}G^\phi(T) + {^{\text{id}}}G^\phi(T) + {^{\text{E}}}G^\phi(T) + {^{\text{mag}}}G^\phi(T)
\end{equation}
\textbf{Implementation}: Evaluated in \texttt{HEAThermoEngine.m} as the thermodynamic objective function minimized during temperature and composition stepping.

\noindent\textbf{Variables}:
\begin{itemize}
  \item $G_m^\phi(T, \mathbf{c})$: Molar Gibbs free energy of phase $\phi$ at absolute temperature $T$ [$\text{J}\cdot\text{mol}^{-1}$].
  \item $T$: Absolute temperature [$\text{K}$].
  \item $\mathbf{c} = (c_1, c_2, \dots, c_n)$: Mole fraction vector of constituent alloy components [$\sum_{i=1}^n c_i = 1$].
  \item ${^{\text{ref}}}G^\phi(T)$: Reference energy of the pure components in the structural state of phase $\phi$ [$\text{J}\cdot\text{mol}^{-1}$].
  \item ${^{\text{id}}}G^\phi(T)$: Ideal mixing contribution [$\text{J}\cdot\text{mol}^{-1}$].
  \item ${^{\text{E}}}G^\phi(T)$: Excess Gibbs free energy accounting for atomic pair interactions [$\text{J}\cdot\text{mol}^{-1}$].
  \item ${^{\text{mag}}}G^\phi(T)$: Magnetic ordering contribution from the Inden-Hillert-Jarl formulation [$\text{J}\cdot\text{mol}^{-1}$].
\end{itemize}

\subsection{Reference Energy from SGTE Unary Lattice Stabilities}
\begin{equation}
{^{\text{ref}}}G^\phi(T) = \sum_{i=1}^n c_i \, {^\circ}G_i^\phi(T)
\end{equation}
\begin{equation}
{^\circ}G_i^\phi(T) - H_i^{\text{SER}} = a_i + b_i T + c_i T \ln T + d_i T^2 + e_i T^{-1} + f_i T^3
\end{equation}
\textbf{Implementation}: Evaluated in \texttt{HEATDBEngine.m} using polynomial coefficients from the SGTE unary database (Dinsdale, 1991).

\noindent\textbf{Variables}:
\begin{itemize}
  \item ${^\circ}G_i^\phi(T)$: Standard molar Gibbs energy of pure element $i$ in phase $\phi$ [$\text{J}\cdot\text{mol}^{-1}$].
  \item $H_i^{\text{SER}}$: Enthalpy of pure element $i$ in its Standard Element Reference (SER) state at $298.15\text{ K}$ and $10^5\text{ Pa}$ [$\text{J}\cdot\text{mol}^{-1}$].
  \item $a_i, b_i, c_i, d_i, e_i, f_i$: Temperature-dependent polynomial coefficients.
  \item $c_i$: Mole fraction of element $i$ [$\text{mol/mol}$].
\end{itemize}

\subsection{Ideal Solution Configurational Entropy}
\begin{equation}
{^{\text{id}}}G^\phi(T) = R T \sum_{i=1}^n c_i \ln c_i
\end{equation}
\textbf{Implementation}: Evaluated in \texttt{HEAThermoEngine.m} for random mixing on the crystal sublattice.

\noindent\textbf{Variables}:
\begin{itemize}
  \item $R$: Universal gas constant, $8.314462 \, \text{J}\cdot\text{mol}^{-1}\cdot\text{K}^{-1}$.
  \item $T$: Absolute temperature [$\text{K}$].
  \item $c_i$: Mole fraction of component $i$ ($0 < c_i \le 1$).
\end{itemize}

\subsection{Redlich-Kister Multicomponent Excess Gibbs Energy}
\begin{equation}
{^{\text{E}}}G^\phi(T) = \sum_{i=1}^{n-1}\sum_{j=i+1}^n c_i c_j \sum_{k=0}^p L_{ij}^{\phi, (k)}(T) \, (c_i - c_j)^k, \quad L_{ij}^{\phi, (k)}(T) = a_k + b_k T
\end{equation}
\textbf{Implementation}: Evaluated in \texttt{HEAThermoEngine.m} using binary and ternary Redlich-Kister interaction coefficients with Muggianu geometric extrapolation.

\noindent\textbf{Variables}:
\begin{itemize}
  \item $L_{ij}^{\phi, (k)}(T)$: $k$-th order binary interaction parameter between elements $i$ and $j$ in phase $\phi$ [$\text{J}\cdot\text{mol}^{-1}$].
  \item $a_k$: Enthalpic interaction coefficient [$\text{J}\cdot\text{mol}^{-1}$].
  \item $b_k$: Entropic interaction coefficient [$\text{J}\cdot\text{mol}^{-1}\cdot\text{K}^{-1}$].
  \item $p$: Expansion order ($p = 0$ for regular solutions; $p \ge 1$ for sub-regular solutions).
\end{itemize}

\subsection{Chemical Potential Balance and Multiphase Equilibrium}
\begin{equation}
\mu_i^\alpha(T, \mathbf{x}^\alpha) = \mu_i^\beta(T, \mathbf{x}^\beta) = \dots = \mu_i^\nu(T, \mathbf{x}^\nu) \quad (\forall i = 1, \dots, n)
\end{equation}
\begin{equation}
\sum_{\phi=1}^\Phi f_\phi x_i^\phi = c_i^0, \quad \sum_{\phi=1}^\Phi f_\phi = 1
\end{equation}
\textbf{Implementation}: Solved in \texttt{HEAThermoEngine.m} using a constrained Newton-Raphson scheme to match chemical potentials under mass conservation.

\noindent\textbf{Variables}:
\begin{itemize}
  \item $\mu_i^\phi$: Chemical potential of component $i$ in phase $\phi$ [$\text{J}\cdot\text{mol}^{-1}$], where $\mu_i^\phi = (\partial G^\phi / \partial n_i)_{T, P, n_{j \ne i}}$.
  \item $f_\phi$: Molar phase fraction of phase $\phi$ ($0 \le f_\phi \le 1$).
  \item $\mathbf{x}^\phi = (x_1^\phi, \dots, x_n^\phi)$: Composition of phase $\phi$ [$\text{mol/mol}$].
  \item $c_i^0$: Overall alloy mole fraction of component $i$.
  \item $\alpha, \beta, \dots, \nu$: Equilibrium phases (e.g. FCC\_A1, BCC\_A2, B2, Liquid, Sigma).
\end{itemize}

\section{Empirical Solid-Solution Criteria and Solidification}
Primary Module: \texttt{HEAThermoEngine.m}

\subsection{Miedema Enthalpy of Mixing and Configurational Entropy}
\begin{equation}
\Delta H_{\text{mix}} = \sum_{i=1}^{n-1}\sum_{j=i+1}^n 4 \, \Omega_{ij} \, c_i \, c_j, \quad \Delta S_{\text{mix}} = -R \sum_{i=1}^n c_i \ln c_i
\end{equation}
\textbf{Implementation}: Evaluated in \texttt{HEAThermoEngine.m} using Miedema regular-solution interaction enthalpies.

\noindent\textbf{Variables}:
\begin{itemize}
  \item $\Delta H_{\text{mix}}$: Molar enthalpy of mixing [$\text{kJ}\cdot\text{mol}^{-1}$].
  \item $\Omega_{ij}$: Equiatomic binary mixing enthalpy between elements $i$ and $j$ [$\text{kJ}\cdot\text{mol}^{-1}$].
  \item $\Delta S_{\text{mix}}$: Configurational entropy of mixing [$\text{J}\cdot\text{mol}^{-1}\cdot\text{K}^{-1}$].
  \item $c_i, c_j$: Atomic fractions of components $i$ and $j$.
\end{itemize}

\subsection{Atomic Size Mismatch ($\delta$) and Electronegativity Difference ($\Delta\chi$)}
\begin{equation}
\delta = \sqrt{\sum_{i=1}^n c_i \left(1 - \frac{r_i}{\bar{r}}\right)^2} \times 100\%, \quad \bar{r} = \sum_{i=1}^n c_i r_i
\end{equation}
\begin{equation}
\Delta\chi = \sqrt{\sum_{i=1}^n c_i \left(\chi_i - \bar{\chi}\right)^2}, \quad \bar{\chi} = \sum_{i=1}^n c_i \chi_i
\end{equation}
\textbf{Implementation}: Evaluated in \texttt{HEAThermoEngine.m} to assess Hume-Rothery topological strain criteria ($\delta \le 6.6\%$) and chemical electronegativity differences.

\noindent\textbf{Variables}:
\begin{itemize}
  \item $\delta$: Atomic size difference parameter [$\%$].
  \item $r_i$: Goldschmidt atomic radius of element $i$ [$\text{\AA}$].
  \item $\bar{r}$: Average atomic radius of the alloy system [$\text{\AA}$].
  \item $\chi_i$: Pauling electronegativity of element $i$.
  \item $\Delta\chi$: Root-mean-square electronegativity mismatch.
\end{itemize}

\subsection{Yang-Zhang Thermodynamic Parameter ($\Omega$)}
\begin{equation}
\Omega = \frac{T_m \, \Delta S_{\text{mix}}}{|\Delta H_{\text{mix}}|}, \quad T_m = \sum_{i=1}^n c_i T_{m,i}
\end{equation}
\textbf{Implementation}: Evaluated in \texttt{HEAThermoEngine.m} using the Yang-Zhang criterion ($\Omega \ge 1.1$ and $\delta \le 6.6\%$).

\noindent\textbf{Variables}:
\begin{itemize}
  \item $\Omega$: Thermodynamic stability parameter [dimensionless].
  \item $T_m$: Rule-of-mixtures melting temperature [$\text{K}$].
  \item $T_{m,i}$: Melting point of pure component $i$ [$\text{K}$].
\end{itemize}

\subsection{Guo Valence Electron Concentration (VEC) Criteria}
\begin{equation}
\text{VEC} = \sum_{i=1}^n c_i \, \text{VEC}_i
\end{equation}
\textbf{Implementation}: Evaluated in \texttt{HEAThermoEngine.m} following the phase boundary criteria of Guo et al.: $\text{VEC} \ge 8.0$ favors FCC; $\text{VEC} < 6.87$ favors BCC; $6.87 \le \text{VEC} < 8.0$ indicates mixed FCC+BCC.

\noindent\textbf{Variables}:
\begin{itemize}
  \item $\text{VEC}$: Valence electron concentration of the alloy [electrons/atom].
  \item $\text{VEC}_i$: Valence electron count of element $i$.
\end{itemize}

\subsection{Scheil-Gulliver Solidification and Homogenization Kinetics}
\begin{equation}
T_{\text{ScheilSolidus}} = T_{\text{Solidus}} - 0.45 \cdot (T_{\text{Liquidus}} - T_{\text{Solidus}}) \cdot (1 + 0.05 \cdot \delta)
\end{equation}
\begin{equation}
I_{\text{seg}} = 0.25 + 0.06 \cdot \delta, \quad T_{\text{homog}} = \min(T_{\text{Solidus}} - 50 \text{ K}, \, 0.88 \cdot T_{\text{Solidus}}), \quad \tau_{\text{homog}} = 10.0 \cdot \left(\frac{I_{\text{seg}}}{0.35}\right)^2 \text{ hours}
\end{equation}
\textbf{Implementation}: Evaluated in \texttt{HEAThermoEngine.m} to estimate solute segregation under rapid cooling and calculate homogenization heat treatment conditions.

\noindent\textbf{Variables}:
\begin{itemize}
  \item $T_{\text{ScheilSolidus}}$: Non-equilibrium solidus temperature [$\text{K}$].
  \item $T_{\text{Solidus}}, T_{\text{Liquidus}}$: Equilibrium solidus and liquidus temperatures [$\text{K}$].
  \item $I_{\text{seg}}$: Segregation severity index.
  \item $T_{\text{homog}}$: Annealing homogenization temperature [$\text{K}$].
  \item $\tau_{\text{homog}}$: Isothermal soaking time [$\text{hours}$].
\end{itemize}

\section{Lattice Distortion and Dislocation Mechanics}
Primary Module: \texttt{HEALatticeDistortionEngine.m}

\subsection{Effective Relaxed Atomic Radius}
\begin{equation}
r_i^* = r_i \left(1 - \beta \frac{K_i - \bar{K}}{\bar{K}}\right), \quad \beta = \frac{1 + \nu}{3(1 - \nu)}, \quad \bar{K} = \sum_{i=1}^n c_i K_i
\end{equation}
\textbf{Implementation}: Evaluated in \texttt{HEALatticeDistortionEngine.m} following the elastic soft-sphere model of He and Yang (2018).

\noindent\textbf{Variables}:
\begin{itemize}
  \item $r_i^*$: Effective relaxed atomic radius of component $i$ in the solid solution [$\text{\AA}$].
  \item $r_i$: Unperturbed metallic radius [$\text{\AA}$].
  \item $K_i, \bar{K}$: Bulk modulus of component $i$ and composition-weighted average bulk modulus [$\text{GPa}$].
  \item $\nu$: Average Poisson's ratio ($\approx 0.30 - 0.33$).
  \item $\beta$: Elastic stress accommodation factor.
\end{itemize}

\subsection{Root-Mean-Square Lattice Strain}
\begin{equation}
\varepsilon_{\text{RMS}} = \sqrt{\sum_{i=1}^n c_i \left(\frac{r_i^* - \bar{r}^*}{\bar{r}^*}\right)^2}, \quad \bar{r}^* = \sum_{i=1}^n c_i r_i^*
\end{equation}
\textbf{Implementation}: Evaluated in \texttt{HEALatticeDistortionEngine.m} to quantify static atomic displacements.

\noindent\textbf{Variables}:
\begin{itemize}
  \item $\varepsilon_{\text{RMS}}$: Root-mean-square static displacement strain [dimensionless].
  \item $\bar{r}^*$: Mean effective atomic radius [$\text{\AA}$].
\end{itemize}

\subsection{Theoretical Shear Strain and Collective Solute Pinning}
\begin{equation}
\bar{\gamma}_{\text{th}} = \gamma_0 \left(1 - \alpha_{\text{dist}} \varepsilon_{\text{RMS}}\right)
\end{equation}
\begin{equation}
\tau_c = \tau_{\text{Peierls}} + \Delta\tau_{\text{Labusch}}, \quad \Delta\tau_{\text{Labusch}} = A_L \cdot G \cdot \varepsilon_{\text{RMS}}^{4/3} \cdot c_i^{2/3}
\end{equation}
\textbf{Implementation}: Evaluated in \texttt{HEALatticeDistortionEngine.m} to compute theoretical shear resistance and the Labusch solute pinning increment.

\noindent\textbf{Variables}:
\begin{itemize}
  \item $\bar{\gamma}_{\text{th}}$: Theoretical plastic shear strain limit.
  \item $\gamma_0$: Ideal shear strain of unperturbed lattice ($\approx 0.11$).
  \item $\alpha_{\text{dist}}$: Lattice distortion coupling coefficient ($\approx 2.8$).
  \item $\tau_c$: Critical resolved shear stress [$\text{MPa}$].
  \item $\tau_{\text{Peierls}}$: Intrinsic lattice friction stress [$\text{MPa}$].
  \item $\Delta\tau_{\text{Labusch}}$: Solid-solution strengthening increment [$\text{MPa}$].
  \item $G$: Shear modulus [$\text{GPa}$].
  \item $A_L$: Dimensionless scaling factor.
\end{itemize}

\subsection{Dislocation Activation Volume}
\begin{equation}
V^* = b \cdot d_w \cdot l_s = -\frac{\partial \Delta G^*}{\partial \tau}
\end{equation}
\textbf{Implementation}: Evaluated in \texttt{HEALatticeDistortionEngine.m} for thermally activated dislocation glide.

\noindent\textbf{Variables}:
\begin{itemize}
  \item $V^*$: Dislocation activation volume [$\text{m}^3$ or in units of $b^3$].
  \item $b$: Magnitude of Burgers vector [$\text{nm}$] ($\approx 0.255\text{ nm}$ for FCC, $\approx 0.286\text{ nm}$ for BCC).
  \item $d_w$: Obstacle barrier width [$\text{nm}$].
  \item $l_s$: Average spacing between pinning obstacles [$\text{nm}$].
\end{itemize}

\section{Mechanical Strengthening and Constitutive Relations}
Primary Module: \texttt{HEAMechanicalEngine.m}

\subsection{Yield Strength Superposition}
\begin{equation}
\sigma_y = \sigma_0 + \Delta\sigma_{\text{ss}} + \Delta\sigma_{\text{HP}} + \Delta\sigma_{\text{ppt}}
\end{equation}
\textbf{Implementation}: Evaluated in \texttt{HEAMechanicalEngine.m} to calculate 0.2\% offset engineering tensile yield strength.

\noindent\textbf{Variables}:
\begin{itemize}
  \item $\sigma_y$: Engineering yield strength [$\text{MPa}$].
  \item $\sigma_0$: Lattice friction stress [$\text{MPa}$] ($\sim 45\text{ MPa}$ for FCC; $\sim 650 - 950\text{ MPa}$ for BCC).
  \item $\Delta\sigma_{\text{ss}}$: Solid-solution strengthening contribution [$\text{MPa}$].
  \item $\Delta\sigma_{\text{HP}}$: Hall-Petch grain boundary strengthening [$\text{MPa}$].
  \item $\Delta\sigma_{\text{ppt}}$: Precipitation strengthening contribution [$\text{MPa}$].
\end{itemize}

\subsection{Solid Solution Strengthening: Varvenne-Curtin (FCC) and Maresca-Curtin (BCC)}
\begin{equation}
\Delta\sigma_{\text{ss}}^{\text{FCC}} = 0.0405 \cdot M \cdot \mu \left(\frac{1 + \nu}{1 - \nu}\right)^{4/3} \left[ \sum_{i=1}^n c_i \left(\frac{\Delta V_i}{V}\right)^2 \right]^{2/3}
\end{equation}
\begin{equation}
\Delta\sigma_{\text{ss}}^{\text{BCC}} = M \cdot \left[ \tau_{\text{kink}} + \alpha_k \cdot \mu \cdot \Delta\varepsilon_{\text{size}}^{4/3} \right]
\end{equation}
\textbf{Implementation}: Evaluated in \texttt{HEAMechanicalEngine.m} based on the random alloy strengthening theories of Varvenne et al. (2016) and Maresca and Curtin (2020).

\noindent\textbf{Variables}:
\begin{itemize}
  \item $M$: Taylor orientation factor ($M = 3.06$ for FCC; $M = 2.73$ for BCC).
  \item $\mu$: Shear modulus [$\text{GPa}$].
  \item $\nu$: Poisson's ratio.
  \item $\Delta V_i / V$: Normalized misfit volume of solute element $i$.
  \item $\tau_{\text{kink}}$: Stress threshold for screw dislocation kink-pair nucleation [$\text{MPa}$].
\end{itemize}

\subsection{Hall-Petch Hardening and High-Temperature Softening}
\begin{equation}
\Delta\sigma_{\text{HP}} = \frac{k_{\text{HP}}}{\sqrt{d}}, \quad \sigma_y(T) = \sigma_y(298\text{ K}) \left[ 1 - \left(\frac{T - 298}{T_m - 298}\right)^m \right]
\end{equation}
\textbf{Implementation}: Evaluated in \texttt{HEAMechanicalEngine.m} for grain-size hardening and thermal softening up to the melting point.

\noindent\textbf{Variables}:
\begin{itemize}
  \item $k_{\text{HP}}$: Hall-Petch coefficient [$\text{MPa}\cdot\mu\text{m}^{1/2}$] ($494 \, \text{MPa}\cdot\mu\text{m}^{1/2}$ for CoCrFeMnNi).
  \item $d$: Mean grain diameter [$\mu\text{m}$] ($30 \, \mu\text{m}$ for standard reference state).
  \item $T_m$: Solidus temperature [$\text{K}$].
  \item $m$: Thermal softening exponent ($0.65 - 0.85$).
\end{itemize}

\subsection{Hardness, True Flow Stress, and Wear Loss}
\begin{equation}
\text{HV} \approx 0.3 \, \sigma_y + 60, \quad \sigma_{\text{true}} = K \, \varepsilon_{\text{true}}^n, \quad W_{\text{wear}} = K_{\text{Archard}} \frac{F_N \, s}{H}, \quad \eta = \frac{A_{\text{sub}}}{A_{\text{clad}} + A_{\text{sub}}} \times 100\%
\end{equation}
\textbf{Implementation}: Evaluated in \texttt{HEAMechanicalEngine.m} for Vickers microhardness, Hollomon work-hardening, Archard wear loss, and cladding dilution.

\noindent\textbf{Variables}:
\begin{itemize}
  \item $\text{HV}$: Vickers microhardness [$\text{kgf}\cdot\text{mm}^{-2}$].
  \item $K$: Hollomon strength coefficient [$\text{MPa}$].
  \item $n$: Strain-hardening exponent.
  \item $W_{\text{wear}}$: Volumetric wear loss [$\text{mm}^3$].
  \item $K_{\text{Archard}}$: Dimensionless Archard wear coefficient.
  \item $F_N$: Normal load [$\text{N}$].
  \item $s$: Sliding distance [$\text{m}$].
  \item $\eta$: Substrate dilution fraction [$\%$].
\end{itemize}

\section{Thermophysical Transport Properties}
Primary Module: \texttt{HEATransportEngine.m}

\subsection{Density, Specific Heat, and Thermal Conductivity}
\begin{equation}
\rho(T) = \frac{\rho(298\text{ K})}{1 + 3 \, \alpha_L (T - 298)}, \quad C_p(T) = \sum_{i=1}^n c_i \, C_{p, i}(T), \quad k(T) = \frac{L_0 \, T}{\rho_{\text{el}}(T)} + \frac{A_{\text{ph}}}{T}
\end{equation}
\textbf{Implementation}: Evaluated in \texttt{HEATransportEngine.m} for temperature-dependent density, Kopp-Neumann heat capacity, and thermal conductivity.

\noindent\textbf{Variables}:
\begin{itemize}
  \item $\rho(T)$: Density at temperature $T$ [$\text{kg}\cdot\text{m}^{-3}$].
  \item $\alpha_L$: Linear coefficient of thermal expansion [$\text{K}^{-1}$].
  \item $C_p(T)$: Molar heat capacity at constant pressure [$\text{J}\cdot\text{mol}^{-1}\cdot\text{K}^{-1}$].
  \item $k(T)$: Total thermal conductivity [$\text{W}\cdot\text{m}^{-1}\cdot\text{K}^{-1}$].
  \item $L_0$: Sommerfeld Lorenz number, $2.443 \times 10^{-8} \, \text{W}\cdot\Omega\cdot\text{K}^{-2}$.
  \item $\rho_{\text{el}}(T)$: Electrical resistivity [$\Omega\cdot\text{m}$].
\end{itemize}

\subsection{Semi-Solid Viscosity and Laser Absorptivity}
\begin{equation}
\eta(f_s) = \eta_{\text{liq}} \left(1 - \frac{f_s}{f_{s, \text{crit}}}\right)^{-2.5 \, f_{s, \text{crit}}}, \quad A_{\text{laser}}(\lambda, T) \approx \sqrt{\frac{8 \pi \varepsilon_0 c_0 \, \rho_{\text{el}}(T)}{\lambda}}
\end{equation}
\textbf{Implementation}: Evaluated in \texttt{HEATransportEngine.m} using the Krieger-Dougherty semi-solid rheology relation and Hagen-Rubens optical absorptivity.

\noindent\textbf{Variables}:
\begin{itemize}
  \item $\eta(f_s)$: Dynamic apparent viscosity in the mushy zone [$\text{Pa}\cdot\text{s}$].
  \item $\eta_{\text{liq}}$: Molten liquid viscosity ($\approx 0.003 - 0.006 \, \text{Pa}\cdot\text{s}$).
  \item $f_{s, \text{crit}}$: Critical solid fraction at rheological jamming ($0.50 - 0.60$).
  \item $A_{\text{laser}}$: Optical surface absorptivity [dimensionless].
  \item $\lambda$: Laser wavelength [$\text{m}$] ($1064\text{ nm}$ for Nd:YAG).
\end{itemize}

\section{Environmental Kinetics: Oxidation and Corrosion}
Primary Modules: \texttt{HEAOxidationEngine.m}, \texttt{HEACorrosionEngine.m}

\subsection{Wagner Parabolic Oxidation and Pilling-Bedworth Ratio}
\begin{equation}
\left(\frac{\Delta W}{A}\right)^2 = k_p(T) \, t, \quad k_p(T) = k_0 \exp\left(-\frac{Q_{\text{ox}}}{R T}\right), \quad \text{PBR} = \frac{M_{\text{ox}} \, \rho_{\text{met}}}{n_{\text{met}} \, M_{\text{met}} \, \rho_{\text{ox}}}
\end{equation}
\textbf{Implementation}: Evaluated in \texttt{HEAOxidationEngine.m} for isothermal oxidation mass gain curves and oxide scale compactness.

\noindent\textbf{Variables}:
\begin{itemize}
  \item $\Delta W / A$: Mass gain per unit area [$\text{mg}\cdot\text{cm}^{-2}$].
  \item $k_p(T)$: Parabolic rate constant [$\text{mg}^2\cdot\text{cm}^{-4}\cdot\text{h}^{-1}$].
  \item $t$: Exposure duration [$\text{hours}$].
  \item $Q_{\text{ox}}$: Oxidation activation energy [$\text{kJ}\cdot\text{mol}^{-1}$].
  \item $\text{PBR}$: Pilling-Bedworth ratio ($1.0 \le \text{PBR} \le 2.0$ indicates protective scale).
\end{itemize}

\subsection{Electrochemical Corrosion Rate (ASTM G102) and PREN}
\begin{equation}
\text{CR} = K_{\text{corr}} \, \frac{i_{\text{corr}} \, \text{EW}}{\rho}, \quad \text{EW} = \left[ \sum_{i=1}^n \frac{w_i \, n_i}{M_i} \right]^{-1}, \quad \text{PREN}_{\text{HEA}} = \text{wt}\%\text{Cr} + 3.3 (\text{wt}\%\text{Mo} + 0.5\,\text{wt}\%\text{W}) + 16\,\text{wt}\%\text{N}
\end{equation}
\begin{equation}
i_{\text{net}} = i_{\text{corr}} \left[ \exp\left(\frac{2.303(E - E_{\text{corr}})}{\beta_a}\right) - \exp\left(-\frac{2.303(E - E_{\text{corr}})}{\beta_c}\right) \right]
\end{equation}
\textbf{Implementation}: Evaluated in \texttt{HEACorrosionEngine.m} according to ASTM G102-89 and Butler-Volmer polarization kinetics.

\noindent\textbf{Variables}:
\begin{itemize}
  \item $\text{CR}$: Annual corrosion penetration rate [$\text{mm}\cdot\text{year}^{-1}$].
  \item $K_{\text{corr}}$: Conversion factor, $0.00327 \, \text{mm}\cdot\text{g}\cdot\mu\text{A}^{-1}\cdot\text{cm}^{-1}\cdot\text{year}^{-1}$.
  \item $i_{\text{corr}}$: Corrosion current density [$\mu\text{A}\cdot\text{cm}^{-2}$].
  \item $\text{EW}$: Alloy equivalent weight [$\text{g}\cdot\text{equivalent}^{-1}$].
  \item $w_i$: Mass fraction of element $i$.
  \item $n_i$: Primary valence state in dissolution (e.g. Fe: 2, Cr: 3, Ni: 2, Mo: 4).
  \item $\text{PREN}_{\text{HEA}}$: Pitting resistance equivalent number.
  \item $\beta_a, \beta_c$: Anodic and cathodic Tafel slopes [$\text{V}\cdot\text{decade}^{-1}$].
\end{itemize}

\section{Kinematic Powder X-Ray Diffractometry}
Primary Module: \texttt{HEAXRDEngine.m}

\subsection{Bragg Diffraction and Structure Factor}
\begin{equation}
\lambda = 2 \, d_{hkl} \, \sin\theta_{hkl}, \quad d_{hkl} = \frac{a_0}{\sqrt{h^2 + k^2 + l^2}}
\end{equation}
\begin{equation}
I_{hkl} \propto |F_{hkl}|^2 \cdot p_{hkl} \cdot \left(\frac{1 + \cos^2(2\theta)}{\sin^2\theta \cos\theta}\right) \cdot \exp(-2 B_{\text{DW}} \sin^2\theta / \lambda^2)
\end{equation}
\textbf{Implementation}: Evaluated in \texttt{HEAXRDEngine.m} to synthesize kinematic X-ray diffraction patterns for cubic and hexagonal phases.

\noindent\textbf{Variables}:
\begin{itemize}
  \item $\lambda$: X-ray wavelength [$\text{\AA}$] ($1.54056 \, \text{\AA}$ for $\text{Cu-K}\alpha_1$).
  \item $\theta_{hkl}$: Bragg angle [$\text{degrees}$].
  \item $d_{hkl}$: Interplanar spacing of $(hkl)$ planes [$\text{\AA}$].
  \item $a_0$: Lattice parameter from Vegard's law with soft-sphere relaxation [$\text{\AA}$].
  \item $F_{hkl}$: Structure factor for reflection $(hkl)$.
  \item $p_{hkl}$: Multiplicity factor of the reflection plane family.
  \item $B_{\text{DW}}$: Debye-Waller parameter [$\text{\AA}^2$].
\end{itemize}

\subsection{Peak Broadening: Scherrer and Williamson-Hall}
\begin{equation}
\text{FWHM}_{hkl} = \frac{K_{\text{Scherrer}} \, \lambda}{D_{\text{crystallite}} \cos\theta_{hkl}} + 4 \, \varepsilon_{\text{micro}} \tan\theta_{hkl}
\end{equation}
\textbf{Implementation}: Evaluated in \texttt{HEAXRDEngine.m} to simulate line broadening from finite crystallite size and lattice microstrain.

\noindent\textbf{Variables}:
\begin{itemize}
  \item $\text{FWHM}_{hkl}$: Peak full width at half maximum [$\text{rad}$].
  \item $K_{\text{Scherrer}}$: Scherrer shape factor ($\approx 0.94$).
  \item $D_{\text{crystallite}}$: Coherent domain size [$\text{nm}$].
  \item $\varepsilon_{\text{micro}}$: Root-mean-square microstrain.
\end{itemize}

\section{Application Ranking and COMSOL Export}
Primary Modules: \texttt{HEAApplicationDatabase.m}, \texttt{CALPHADHEACalc.m}

\subsection{Multi-Criteria Suitability Scoring}
\begin{equation}
S_j = \left[ \sum_{k=1}^m w_k \cdot \phi_k(P_k) \right] \times \prod_{l=1}^p \Theta_l, \quad \phi_k(P_k) = \frac{1}{1 + \exp\left(-\kappa_k \cdot \frac{P_k - P_{k, \text{target}}}{\Delta P_k}\right)}
\end{equation}
\textbf{Implementation}: Evaluated in \texttt{HEAApplicationDatabase.m} to rank candidate alloys across 8 engineering application sectors.

\noindent\textbf{Variables}:
\begin{itemize}
  \item $S_j$: Overall suitability score for sector $j$ ($0 \le S_j \le 1.0$).
  \item $w_k$: Weighting factor for property $k$ ($\sum w_k = 1.0$).
  \item $\phi_k(P_k)$: Sigmoidal response function for property $P_k \in [0, 1]$.
  \item $\Theta_l \in \{0, 1\}$: Hard constraint flag (e.g. elongation $< 5\%$ disqualifies cryogenic applications).
\end{itemize}

\subsection{Continuous Piecewise Property Smoothing for COMSOL}
\begin{equation}
P_{\text{smooth}}(T) = P_{\text{solid}}(T) + \left[P_{\text{liquid}}(T) - P_{\text{solid}}(T)\right] \cdot \text{flc2hs}\left(T - T_{\text{Solidus}}, \, \Delta T_{\text{smooth}}\right)
\end{equation}
\textbf{Implementation}: Evaluated in \texttt{CALPHADHEACalc.m} to export continuous material property functions for finite element simulations.

\noindent\textbf{Variables}:
\begin{itemize}
  \item $P_{\text{smooth}}(T)$: Continuous property function across solid, mushy, and liquid regimes.
  \item $\text{flc2hs}(x, \Delta x)$: COMSOL smoothed Heaviside function with continuous second derivative.
  \item $\Delta T_{\text{smooth}}$: Smoothing transition temperature width [$\text{K}$].
\end{itemize}

\section*{References}
\begin{enumerate}
  \item A. T. Dinsdale, ``SGTE data for pure elements,'' \textit{CALPHAD}, vol. 15, no. 4, pp. 317--425, 1991. DOI: 10.1016/0364-5916(91)90030-N.
  \item Q. He and Y. Yang, ``On lattice distortion in high entropy alloys,'' \textit{Frontiers in Materials}, vol. 5, p. 42, 2018. DOI: 10.3389/fmats.2018.00042.
  \item C. Varvenne, A. Luque, and W. A. Curtin, ``Theory of strengthening in fcc high entropy alloys,'' \textit{Acta Materialia}, vol. 118, pp. 164--176, 2016. DOI: 10.1016/j.actamat.2016.07.040.
  \item F. Maresca and W. A. Curtin, ``Theory of screw dislocation glide in BCC high-entropy alloys,'' \textit{Acta Materialia}, vol. 182, pp. 235--249, 2020. DOI: 10.1016/j.actamat.2019.10.015.
  \item S. Guo, C. Ng, J. Lu, and C. T. Liu, ``Effect of valence electron concentration on stability of fcc or bcc phase in high entropy alloys,'' \textit{Journal of Applied Physics}, vol. 109, no. 10, p. 103505, 2011. DOI: 10.1063/1.3587228.
  \item ASTM G102-89(2015)e1, ``Standard Practice for Calculation of Corrosion Rates and Related Information from Electrochemical Measurements,'' \textit{ASTM International}, West Conshohocken, PA, 2015. DOI: 10.1520/G0102-89R15E01.
  \item A. R. Miedema, P. F. de Ch{\^a}tel, and F. R. de Boer, ``Cohesion in alloys---fundamentals of a semi-empirical model,'' \textit{Physica B+C}, vol. 100, no. 1, pp. 1--28, 1980. DOI: 10.1016/0378-4363(80)90054-6.
  \item X. Yang and Y. Zhang, ``Prediction of high-entropy alloys with a new thermodynamical parameter,'' \textit{Materials Chemistry and Physics}, vol. 132, no. 2--3, pp. 233--238, 2012. DOI: 10.1016/j.matchemphys.2011.11.021.
  \item O. Redlich and A. T. Kister, ``Algebraic representation of thermodynamic properties and the classification of solutions,'' \textit{Industrial \& Engineering Chemistry}, vol. 40, no. 2, pp. 345--348, 1948. DOI: 10.1021/ie50458a035.
  \item M. Hillert and M. Jarl, ``A model for alloying effects in ferromagnetic metals,'' \textit{CALPHAD}, vol. 2, no. 3, pp. 227--238, 1978. DOI: 10.1016/0364-5916(78)90011-1.
  \item E. Scheil, ``Bemerkungen zur Kristallisation von Legierungen,'' \textit{Zeitschrift f{\"u}r Metallkunde}, vol. 34, pp. 70--72, 1942.
  \item J. F. Archard, ``Contact and rubbing of flat surfaces,'' \textit{Journal of Applied Physics}, vol. 24, no. 8, pp. 981--988, 1953. DOI: 10.1063/1.1721448.
  \item C. Wagner, ``Beitrag zur Theorie des Anlaufvorgangs,'' \textit{Zeitschrift f{\"u}r Physikalische Chemie}, vol. 21B, no. 1, pp. 25--41, 1933.
  \item P. Scherrer, ``Bestimmung der Gr{\"o}sse und der inneren Struktur von Kolloidteilchen mittels R{\"o}ntgenstrahlen,'' \textit{Nachrichten von der Gesellschaft der Wissenschaften zu G{\"o}ttingen}, vol. 1918, pp. 98--100, 1918.
\end{enumerate}

\end{document}
